\documentclass[10pt]{article}
\usepackage[T1]{fontenc}
\usepackage[utf8]{inputenc}
\usepackage[french]{babel}
\usepackage{amsmath,amssymb,mathrsfs,mathpazo}
\usepackage{geometry}
\geometry{a4paper,margin=2cm}
\usepackage{graphicx,array,tabularx,tabularray,colortbl,multirow,multicol,booktabs,enumitem,pifont,ulem,fancybox,diagbox,cancel}
\usepackage{pstricks,pstricks-add,pst-plot,pst-tree,pst-node,pst-text,pst-func,pst-math,pst-bspline,pst-3dplot,pst-eucl,pst-solides3d}
\let\pscalculate\undefined
\usepackage{pst-bezier}
\usepackage{tikz,pgfplots,tkz-tab}
\usetikzlibrary{arrows,arrows.meta,patterns,calc,positioning,trees,decorations.markings,shapes.geometric,shapes.misc,shapes.arrows,intersections,decorations.pathmorphing,decorations.pathreplacing,backgrounds,shadows,babel}
\pgfplotsset{compat=1.18}
\usepgfplotslibrary{fillbetween}
\usepackage[np]{numprint}
\usepackage{listings,makecell,pas-tableur,twoopt}
\let\barre\undefined
\usepackage[pstricks]{tablvar}
\definecolor{darkgreen}{RGB}{0,100,0}
\providecommand{\red}{\color{red}}
\providecommand{\blue}{\color{blue}}
\providecommand{\green}{\color{green}}
\providecommand{\black}{\color{black}}
\providecommand{\white}{\color{white}}
\providecommand{\gray}{\color{gray}}
\providecommand{\cyan}{\color{cyan}}
\providecommand{\magenta}{\color{magenta}}
\providecommand{\yellow}{\color{yellow}}
\providecommand{\violet}{\color{violet}}
\providecommand{\orange}{\color{orange}}
\providecommand{\euro}{\texteuro}
\providecommand{\EUR}{\texteuro}
\providecommand{\ieme}{\textsuperscript{e}}
\providecommand{\degres}{\ensuremath{^\circ}}
\providecommand{\degree}{\ensuremath{^\circ}}
\providecommand{\vv}[1]{\overrightarrow{#1}}
\providecommand{\Coord}[1]{#1}
\providecommand{\decosix}{}
\providecommand{\starredbullet}{\ensuremath{\bullet}}
\pagestyle{empty}
\setlength{\parindent}{0pt}
\usepackage{adjustbox}\DeclareUnicodeCharacter{2236}{:}\DeclareUnicodeCharacter{2190}{\ensuremath{\leftarrow}}
\providecommand{\R}{\mathbb{R}}
\providecommand{\N}{\mathbb{N}}
\providecommand{\Z}{\mathbb{Z}}
\providecommand{\Q}{\mathbb{Q}}
\providecommand{\D}{\mathbb{D}}
\providecommand{\C}{\mathbb{C}}
\providecommand{\e}{{\mathrm{e}}}
\providecommand{\vect}[1]{\overrightarrow{#1}}
\providecommand{\vectt}[1]{\overrightarrow{\mathrm{#1}}}
\providecommand{\vvt}[1]{\overrightarrow{\mathrm{#1}}}
\providecommand{\np}[2][]{{#2~#1}}
\providecommand{\vv}[1]{{\overrightarrow{#1}}}
\renewcommand{\d}{{\mathrm{d}}}
\providecommand{\degres}{{^\circ}}
\providecommand{\degre}{{^\circ}}
\providecommand{\blue}{\color{blue}}
\providecommand{\red}{\color{red}}
\providecommand{\enskip}{\quad}
\providecommand{\llbracket}{\lbrack\!\lbrack}
\providecommand{\rrbracket}{\rbrack\!\rbrack}
\providecommand{\euro}{{\text{€}}}
\providecommand{\no}{{\text{n°}}}
\providecommand{\interval}[2]{{[#1;#2]}}
\providecommand{\intervFO}[2]{{[#1;#2[}}
\providecommand{\dotfill}{\ldots}

\providecommand{\ds}{\displaystyle}
\providecommand{\textding}[1]{#1}
\providecommand{\Oij}{\(\mathrm{O};\vec{\imath},\vec{\jmath}\)}
\providecommand{\Oijk}{\(\mathrm{O};\vec{\imath},\vec{\jmath},\vec{k}\)}
\providecommand{\Ouv}{\(\mathrm{O};\vec{u},\vec{v}\)}
\providecommand{\intervFF}[2]{\ensuremath{[#1;#2]}}
\renewcommand{\euro}{\ifmmode\text{\texteuro}\else\texteuro\fi}
\DeclareUnicodeCharacter{00B0}{\ifmmode^{\circ}\else\textdegree\fi}
\DeclareUnicodeCharacter{20AC}{\ifmmode\text{\texteuro}\else\texteuro\fi}
% Énoncés : Éducation nationale. Transcription et corrigés : APMEP, Denis Vergès et contributeurs. Découpage et classement : CoopMaths, Sébastien Lozano, Éric Elter et contributeurs. Adaptation de la présentation : Maths Explorer.
% CC BY-SA : https://coopmaths.fr/www/about/
% https://coopmaths.fr/alea/static/bac/2024/tex/bac_2024_01_sujet0_7.tex
\begin{document}
\noindent\textit{APMEP / CoopMaths. CC BY-SA. Présentation adaptée par Maths Explorer.}\par\medskip
On considère les fonctions $f_{k}$ définies sur $\R$ par 
\[f_{k}(x) = x + k \e^{-x},\]
 où $k$ est un réel strictement positif.

	\begin{enumerate}
		\item On s'intéresse dans cette question au cas $k = 0,5$, donc à la fonction $f_{0,5}$ définie sur $\R$ par :
\[f_{0,5}(x) = x + 0,5 \e^{-x}.\]

\smallskip
	\begin{enumerate}
		\item Montrer que la dérivée de $f_{0,5}$, notée $f'{ }_{0,5}$ vérifie $f'{ }_{0,5}(x) = 1-0,5 \e^{-x}$.

		\item Montrer que la fonction $f_{0,5}$ admet un minimum en $\ln (0,5)$.
	\end{enumerate}
\end{enumerate}

		Soit $k$ un réel strictement positif. On donne le tableau de variations de la fonction $f_{k}$.
		\medskip
		

					\hfill~\begin{tikzpicture}
				\tkzTabInit[lgt=3]{Valeurs de $x$/0.7,
					Variations de $f_k$/1.5}{$-\infty$,$\ln(k)$,$+\infty$}
				\tkzTabVar{+/$+ \infty$,-/$f_k(\ln k)$,+/$+\infty$}
			\end{tikzpicture}\hfill~

	\begin{enumerate}
		\setcounter{enumi}{1}
		\item Montrer que pour tout réel positif $k, f_{k}(\ln k) = \ln k + 1$.
	\end{enumerate}

On note $\mathcal{C}_{k}$ la courbe représentative de la fonction $f_{k}$ dans un plan muni d'un repère orthonormé.

On note $\mathrm{A}_{k}$ le point de la courbe $\mathcal{C}_{k}$ d'abscisse $\ln k$.

On a représenté ci-dessous quelques courbes $\mathcal{C}_{k}$ pour différentes valeurs de $k$.

	\begin{center}
			\begin{center}
				\begin{tikzpicture}
					\begin{axis}[x={10mm},y={10mm},    xmin=-4,  xmax=8.9,    ymin = 0,     ymax=8.9, xtick distance=1,    ytick distance = 1,
						tick label style={font=\footnotesize}, minor tick num = 0, grid=none,  axis y line=center,axis x line=bottom]
						\foreach \coef in{0.5,1,1.5,2,2.5,3,3.5,4,4.5,5}
						{\addplot [line width=1.2pt, color=red, smooth, samples=100, domain= -4:9 ]{x+\coef*exp(-x)};}
				\node[label={90:{$\mathrm{A}_{0,5}$}},circle,fill,inner sep=2pt] at (axis cs:-0.693,0.307) {};
				\node[label={0:{$\mathrm{A}_1$}},circle,fill,inner sep=2pt] at (axis cs:0,1) {};
				\node[label={90:{$\mathrm{A}_{1,5}$}},circle,fill,inner sep=2pt] at (axis cs:0.325,1.405) {};
				\node[circle,fill,inner sep=2pt] at (axis cs:0.693,1.693) {};
				\node[circle,fill,inner sep=2pt] at (axis cs:0.916,1.916) {};
				\node[circle,fill,inner sep=2pt] at (axis cs:1.099,2.099) {};
				\node[circle,fill,inner sep=2pt] at (axis cs:1.253,2.253) {};
				\node[circle,fill,inner sep=2pt] at (axis cs:1.386,2.386) {};
				\node[circle,fill,inner sep=2pt] at (axis cs:1.504,2.504) {};
				\node[label={90:{$\mathrm{A}_5$}},circle,fill,inner sep=2pt] at (axis cs:1.609,2.609) {};
			\end{axis}
				\end{tikzpicture}
			\end{center}
		\end{center}

\begin{enumerate}
		\setcounter{enumi}{2}
		\item Indiquer si l'affirmation suivante est vraie ou fausse. Justifier la réponse. Une réponse non justifiée ne rapporte aucun point.

		\begin{itemize}
		\item \textbf{Affirmation :} \og{}Pour tout réel $k$ strictement positif, les points $A_{0,5}, A_{1}$ et $A_{k}$ sont alignés. \fg{}
		\end{itemize}
\end{enumerate}

\end{document}
